\section{Necessary equations at a three-coordinate threshold}
\label{sec:triangular-contact}

We derive necessary conditions for a maximum of the ratio over all
independent arrays when three coordinates have the threshold pattern
described below.

Let the first coordinate have three ordered states, and let the other
two coordinates have two states each. Number states from zero. Suppose
the configurations attaining a selected closed threshold are exactly
\begin{equation}\label{eq:triangle-configurations}
                    (0,0,1),\qquad(0,1,0),\qquad(1,0,0).
\end{equation}
The configuration $(0,0,0)$ lies strictly below the threshold and all
other configurations lie strictly above it. Equality of the three sums
in \eqref{eq:triangle-configurations} forces the two two-point diameters
to equal the lower gap of the three-point law. The upper gap of that
law is unrestricted. Thus the common-diameter theorem does not apply.

\begin{proposition}[Equal probabilities at a triangular threshold]
\label{prop:triangle-equal-probabilities}
At an attained ratio maximum $C>0$ with the configuration
\eqref{eq:triangle-configurations}, the two two-point coordinates have
equal upper-state probabilities.
\end{proposition}
\begin{proof}
Write their common diameter as $d$, their upper probabilities as
$p_1,p_2$, and the triple gaps as $u=d,v>0$. Its signed probability
measure
\[
 \eta_X=\frac1{uv}\left\{
 \frac{v}{u+v}\delta_{x_0}-\delta_{x_1}
                  +\frac{u}{u+v}\delta_{x_2}\right\}
\]
has mass and first moment zero and second moment one. Hold
$s=p_1+p_2$ fixed. If $p_1\ne p_2$, their product $t=p_1p_2$ is a
two-sided local parameter, and the raw pair-sum probabilities are
$(1-s+t,s-2t,t)$ at $0,d,2d$. Its derivative per unit variance is
$\eta_B=(\delta_0-2\delta_d+\delta_{2d})/(2d^2)$.
Transfer variance $\lambda$ to the triple and $-\lambda$ to the
independent pair. The mean and total variance $V$ are fixed, and the
three original laws remain probability laws for both signs of small
$\lambda$. For the threshold in \eqref{eq:triangle-configurations},
the mixed event coefficient is
\[
 g=-\frac{1+v/(u+v)}{2d^2uv}<0,
 \qquad F(\lambda)=F(0)+L\lambda-g\lambda^2.
\]
Indeed the second differences of its three rows at pair sums $0,d,2d$
are $-1,1,0$.

Put $W=p_1(1-p_1)+p_2(1-p_2)$. At fixed $s$, expansion of the two
original third absolute moments gives
\[
 d^{-3}\sum_{j=1}^2\E|Y_j|^3
             =-W^2+(3-4s+2s^2)W+\text{constant}(s).
\]
Their second derivative with respect to pair variance is $-2/d$;
the triple's third absolute moment is affine. The selected deficit
$CR(\lambda)V^{-3/2}+\Phi(y/\sqrt V)-F(\lambda)$ therefore has
second derivative $2g-2C/(dV^{3/2})<0$. It cannot have a local minimum
at the attained law. Thus $p_1=p_2$.
\end{proof}

Normalize the original total variance to one, write its third-moment
sum as $R$, and let $z$ be the selected threshold. For a hypothetical
maximum $C>\be$, put
\[
 K=R+\frac{z\phi(z)}{3C}>0.
\]
Positivity follows from $R\ge1/\sqrt3$, $|z\phi(z)|<1/4$, and
$C>2/5$. Measure the three original coordinates in units of $K$:
\begin{equation}\label{eq:triangle-normalized-law}
 X/K\in\{-a,c,b\},\quad \Pp(X/K=-a,c,b)=(l,m,r),
 \qquad Y_j/K=g(B_p-p),\quad j=1,2,
\end{equation}
where $a,b>0$, $-a<c<b$, $l,m,r>0$, $q=1-p$, and
$g=a+c>0$. Put $v=b-c>0$, $h=a+b=g+v$, and $\rho=p^2+q^2$.
The letters $Q,B_3,y$ below denote total variance, additive third moment,
and threshold in these units; they are not variance-one quantities.

\begin{theorem}[Necessary contact equations for the triangular threshold]
\label{thm:triangle-projection}
Suppose \eqref{eq:triangle-normalized-law} attains the full independent
ratio maximum $C>\be$ at \eqref{eq:triangle-configurations}. Define
\begin{align}
 P_b(x)&=|x|^3-\tfrac32x^2-3b(b-1)x+b^2(2b-\tfrac32),
                                                        \label{eq:triangle-polynomial}\\
 A_0&=P_b(-a),\quad A_1=P_b(c),\quad
 H_0=3\{b(b-1)+a(a-1)\},\quad
 H_1=3\{b(b-1)+c-c|c|\}.\notag
\end{align}
Then $b\ge1$, $A_0>A_1>0$, $H_0,H_1\ge0$, and the probabilities
and contact multiplier are determined by the support geometry:
\begin{equation}\label{eq:triangle-heights}
 p=\frac{A_0-A_1}{A_0+A_1},\qquad
 q=\frac{2A_1}{A_0+A_1},\qquad
 \tau=\frac{4A_1}{(A_0+A_1)^2}=CK^3.
\end{equation}
For fixed $a,c,b$, define the following affine functions of $Q$:
\begin{align}
 m(Q)&=\frac{ab+2pqg^2-Q}{gv},&
 l(Q)&=\frac{b-m(Q)v}{h},&
 r(Q)&=\frac{a-m(Q)g}{h},\label{eq:triangle-affine-masses}\\
 \sigma(Q)&=l(Q)H_0+m(Q)H_1,\notag\\
 B_3(Q)&=l(Q)a^3+m(Q)|c|^3+r(Q)b^3+2pq\rho g^3,
 &y&=c-2pg.\notag
\end{align}
The original variance $Q>0$ satisfies
\begin{align}
 B_3(Q)+y\sigma(Q)/3&=Q,                       \label{eq:triangle-normalization}\\
 \tau\sigma(Q)\sqrt Q&=\phi(y/\sqrt Q),        \label{eq:triangle-density}\\
 l(Q)(1-p^2)+m(Q)q^2-\Phi(y/\sqrt Q)&=\tau B_3(Q),
                                                        \label{eq:triangle-ratio}\\
 l(Q)p+m(Q)q&=\tau\{\sigma(Q)g+\tfrac32(1-2p)g^2
                                      -(1-2p)^3g^3\}. \label{eq:triangle-Bernoulli-contact}
\end{align}
All three masses and $\sigma(Q)$ must be strictly positive, and
$C=\tau Q^{3/2}>\be$.

There is a further quadratic inequality in this same original variance.
Set
\begin{align*}
 \gamma&=h+g-3a+\frac{2a^3}{gh}
                         -\frac{2(-c)_+^3}{gv},\\
 \gamma'&=1-\frac{2a^3}{gh^2}
                         +\frac{2(-c)_+^3}{gv^2},\qquad
 T=\frac{3h^2v^2(\gamma')^2}{2b-1}.
\end{align*}
Then
\begin{equation}\label{eq:triangle-Schur}
 (12\gamma-15)Q^2+y(y^2+2Q)\sigma(Q)+TQr(Q)\le0.
\end{equation}
Equations \eqref{eq:triangle-affine-masses}--\eqref{eq:triangle-Schur}
retain both signs of $c$, including $c=0$.
\end{theorem}

\begin{proof}
The contamination argument of Lemma~\ref{lem:full-law-contacts} applies
to each coordinate with the other two held fixed. After division by
$CK^3$, the triple's piecewise cubic majorant has quadratic coefficient $-3/2$.
Its highest contact is unselected and has deletion value zero, so it
is a smooth minimum. Its value and derivative determine
\eqref{eq:triangle-polynomial}. Global nonnegativity at $-b$ gives
$P_b(-b)=6b^2(b-1)\ge0$, hence $b\ge1$.
At the two selected contacts monotonicity from the left gives
$P_b'(-a),P_b'(c)\le0$. Their actual deletion values are
$1-p^2,q^2$. Solving $\tau A_0=1-p^2$, $\tau A_1=q^2$ proves
\eqref{eq:triangle-heights}.

Centering and variance give $l=(b-mv)/h$, $r=(a-mg)/h$, and
$Q=ab-mgv+2pqg^2$, proving \eqref{eq:triangle-affine-masses}.
The linear coefficient of the marginal contact majorant gives
$\sigma=\phi(z)/(CK^2)=lH_0+mH_1$. Its mean identity, the definition
of $K$, and the actual event probability prove
\eqref{eq:triangle-normalization}--\eqref{eq:triangle-ratio}.
The two deletion values for either Bernoulli coordinate are
$l+mq,lq$. Their difference, inserted in the same contact identity,
is \eqref{eq:triangle-Bernoulli-contact}.

For clarity, the affine equation \eqref{eq:triangle-normalization}
is also the support stationarity equation for the common lower gap;
it is not an extra independent condition. Vary the raw lower gap $g$
of the triple and both Bernoullis, keeping the triple's raw outer
span $h$ and all masses fixed. The three threshold configurations
remain equal and all other comparisons remain strict. Differentiation
and recentering give
\[
 lH_0=2p\{\sigma+3qg(1-\rho g)\}.
\]
Direct substitution of the original moments gives exactly
\[
 B_3-Q+y\sigma/3
 =\frac g3\bigl[lH_0-2p\sigma-6pqg(1-\rho g)\bigr].
\]

We derive \eqref{eq:triangle-Schur} from a two-dimensional family of
original laws. Work in the fixed contact units. Keep the triple's
raw lower supports $0,g$, its raw mean $a$, and both Bernoulli laws
fixed. Vary its highest raw support $H$ and the total variance $V$.
With $V_B=2pqg^2$ set
\[
 A(V)=V-V_B+a^2-ga,\quad
 r(V,H)=\frac{A(V)}{H(H-g)},\quad
 m(V,H)=\frac{a-Hr(V,H)}g,\quad l=1-m-r.
\]
These probabilities are positive in a two-sided neighborhood of
$(V,H)=(Q,h)$. The centered threshold $y$ is fixed. Its deletion
values $d_0=1-p^2,d_1=q^2,d_2=0$ remain fixed because the moving
highest state occurs in no selected configuration. The actual event
and third-moment sum have the forms
\[
 F(V,H)=F_0+A(V)f(H),\quad R(V,H)=R_0+A(V)\gamma(H),
\]
where
\[
 f(H)=\frac{d_0-d_2}{gH}+\frac{d_2-d_1}{g(H-g)},\quad
 \gamma(H)=H+g-3a+\frac{2a^3}{gH}
                       -\frac{2(a-g)_+^3}{g(H-g)}.
\]
Here $F_0$ and $R_0$ are independent of $V,H$. The objective
\[
 J(V,H)=F(V,H)-\Phi(y/\sqrt V)-CR(V,H)V^{-3/2}
\]
has zero gradient and negative semidefinite Hessian at the original
law. Differentiation, using $C/Q^{3/2}=\tau$,
$\phi(y/\sqrt Q)=\tau\sigma\sqrt Q$, and
$B_3=Q-y\sigma/3$, gives
\[
 J_{VV}=\frac{\tau}{4Q^3}
       \{(12\gamma-15)Q^2+\sigma y(y^2+2Q)\},\quad
 J_{VH}=\frac{3\tau A(Q)\gamma'}{2Q},\quad
 J_{HH}=-3\tau r(Q)(2b-1)<0.
\]
In computing $J_{VH}$, differentiate both $R(V,H)$ and
$V^{-3/2}$. The last entry also
follows by evaluating the contact majorant on the moving highest
state: its gap has value and derivative zero and second derivative
$3\tau(2b-1)$. Since $A(Q)=r(Q)hv$, the Schur inequality
$J_{VV}-J_{VH}^2/J_{HH}\le0$ is exactly
\eqref{eq:triangle-Schur}.
\end{proof}

\subsection{Finite scalar tests and their boundary cases}
\label{subsec:triangle-interval}
Write
$\sigma(Q)=s_1Q+s_0$, $B_3(Q)=b_1Q+b_0$. The normalization is
\[
              n_1Q+n_0=0,\qquad
 n_1=b_1+ys_1/3-1,\quad n_0=b_0+ys_0/3.
\]
If $n_1\ne0$, only $Q=-n_0/n_1$ is possible. If $n_1=0\ne n_0$,
the geometry is impossible. If both vanish, retain the full interval
of possible original variances.

For any proved upper bound $C_+$ on the independent constant, this
interval is the intersection of
\[
 ((\be/\tau)^{2/3},(C_+/\tau)^{2/3}]
 \quad\text{with}\quad
 \{Q:l(Q)>0, m(Q)>0, r(Q)>0, \sigma(Q)>0\}.
\]
All mass and slope conditions are open. On this interval the density
equation is equivalent, with its sign retained, to
\begin{equation}\label{eq:triangle-density-squared}
       2\pi\tau^2 Q(s_1Q+s_0)^2 e^{y^2/Q}=1.
\end{equation}
The logarithmic derivative of its left side has positive denominator
$Q^2(s_1Q+s_0)$ and numerator
\[
          3s_1Q^2+(s_0-s_1y^2)Q-s_0y^2.
\]
Thus endpoints and at most two interior critical points determine its
range on each interval. The Schur expression
\eqref{eq:triangle-Schur} is a quadratic polynomial in $Q$ and cuts
the interval into at most two intervals, possibly with isolated points.
Linear and constant polynomials are retained, as are a repeated root
and an identically zero polynomial. Endpoint values in the density
test are one-sided limits; equality only at an excluded endpoint is
not admissible. An isolated retained root is tested directly.

These tests can exclude a support geometry before the remaining
contact equations are solved.
